\subsection{Dimension}

THESE PROBLEMS GOT A PUNCH TO EM!!! Well, mainly I struggled on 8. The rest were fine.

\begin{enumerate}
    \item Show that the subspaces of $\mathbb{R}^2$ are precisely $\{0\}$, all lines in $\mathbb{R}^2$ containing the origin, and $\mathbb{R}^2$.
    

    \begin{proof}
        We can split this up by the number of dimensions that the basis of the subspace contains. 

        \begin{enumerate}
            \item 0 dimensions. The only 0 dimensional vector is $\{0\}$.
            \item 1 dimension. This is a single vector, and it will always produce a line after scaling in $\R^2$. 
            \item 2 dimensions. Because the 2 dimensions are linearly independent, we know that it spans all of $\R^2$.
        \end{enumerate}
    \end{proof}


    \item[3] 
        \begin{enumerate}
            \item Let $U = \{p \in \mathcal{P}_4(\mathbb{F}) : p(6) = 0\}$. Find a basis of $U$.
            
            \begin{proof}
                If $p(6) = 0$, then all polynomials in $U$ will have a factor of $(x-6)$ (including 0 by technicality). Thus, we can choose the basis 

                \begin{align*}
                    (x - 6) \\
                    (x - 6)(x) \\
                    (x - 6)(x^2)
                    (x-6)(x^3)
                \end{align*}
            \end{proof}
            \item Extend the basis in (a) to a basis of $\mathcal{P}_4(\mathbb{F})$.
            
            \begin{proof}
                We can simply add the constant, 1. This is now a linearly independent list which has length $\dim \mathcal P_4(\F)$, so it is a basis.
            \end{proof}
            \item Find a subspace $W$ of $\mathcal{P}_4(\mathbb{F})$ such that $\mathcal{P}_4(\mathbb{F}) = U \oplus W$.
            
            \begin{proof}
                It is just the subspace given by the span of the additional vector, in this case 

                $$W = \{a : a \in \F\}$$
            \end{proof}
        \end{enumerate}


        \item[5]
            \begin{enumerate}
                \item Let $U = \{p \in \mathcal{P}_4(\mathbb{F}) : p(2) = p(5)\}$. Find a basis of $U$.
                
                \begin{proof}
                    We have a fixed point, which should reduce the dimension by 1 but not in an entirely clear way\dots

                    We know there must not be a polynomial of degree exactly 1 because that would make it impossible to satisfy the constraints, so we can simply choose polynomials that satisfy the property. Such as\dots

                    \begin{align*}
                        1 \\
                        (x-2)(x-5) \\
                        (x-2)(x-5)x \\
                        (x-2)(x-5)x^2
                    \end{align*}

                    These 3 vectors are clearly linearly independent and are in $U$. We know that it spans $U$ because it has less than 4 dimensions (otherwise it would equal to $\mathcal P_4(\F)$) and since this is a linearly independent list of 3 it can't be any lower.
                \end{proof}
                \item Extend the basis in (a) to a basis of $\mathcal{P}_4(\mathbb{F})$.
                
                \begin{proof}
                    We can plop in $x$ now.
                \end{proof}
                \item Find a subspace $W$ of $\mathcal{P}_4(\mathbb{F})$ such that $\mathcal{P}_4(\mathbb{F}) = U \oplus W$.
                
                \begin{proof}
                    That would be $W = \{ax : a \in \F\}$.
                \end{proof}
            \end{enumerate}


            \item[7]
            \begin{enumerate}
                \item Let $U = \{p \in \mathcal{P}_4(\mathbb{R}) : \int_{-1}^{1} p = 0\}$. Find a basis of $U$.
                
                \begin{proof}
                    We can create a list of 4 linearly independent vectors. 3 is the maximum basis size for a (strict) subspace of $\mathcal P_4(\F)$ so we are sure it spans the whole subspace.

                    \begin{align*}
                        x \\
                        x^3 \\
                        x^2 - \frac13 \\
                        x^4 - \frac15
                    \end{align*}
                \end{proof}
                \item Extend the basis in (a) to a basis of $\mathcal{P}_4(\mathbb{R})$.
                
                \begin{proof}
                    We can add just one more linearly independent element, say, $1$.
                \end{proof}
                \item Find a subspace $W$ of $\mathcal{P}_4(\mathbb{R})$ such that $\mathcal{P}_4(\mathbb{R}) = U \oplus W$.
                
                \begin{proof}
                    \[W = \{a : a \in \F\}\]
                \end{proof}
            \end{enumerate}

            \item[8] Suppose $v_1, \ldots, v_m$ is linearly independent in $V$ and $w \in V$. Prove that
                $$\dim \text{span}(v_1 + w, \ldots, v_m + w) \geq m - 1$$
                
                \begin{proof}
                    We can subtract the first vector from the rest to obtain that 

                    $$v_2 - v_1, \dots, v_m - v_1 \in \linspan(v_1 + w, \dots, v_m + w)$$

                    Therefore, 

                    $$\linspan(v_2 - v_1, \dots, v_m - v_1) \leq \linspan(v_1 + w, \dots, v_m + w)$$

                    Because 

                    $$\linspan(v_1, \dots, v_m) = \linspan(v_1, v_2 - v_1, \dots, v_m - v_1) = m$$

                    We have that 

                    $$\linspan(v_1 + w, \dots, v_m + w) \geq \linspan(v_2 - v_1, \dots, v_m - v_1) = m-1$$
                \end{proof}


            \item[9] Suppose $m$ is a positive integer and $p_0, p_1, \dots, p_m \in \mathcal P(\mathbb F)$ are such that each $p_k$ has degree $k$. Prove that
            

            \[p_0, p_1, \dots, p_m\]

            is a basis of \(\mathcal P_m(\mathbb F)\).

            \begin{proof}

                We prove it by induction. First, $\mathcal P_0(\mathbb F)$ is just $\F$. Any non-zero value in $\F$ is a basis for $\F$. 

                Suppose that $p_0, \dots, p_i$ was a basis for $\mathcal P_i(\F)$. Then, we introduce arbitrary $i + 1$th degree polynomial $p_{i + 1}$. Given some arbitrary polynomial in $\mathcal P_{i + 1}(\F)$, $f(x)$, we can set the coefficient of $p_{i + 1}$ to be the coefficient of the $x^{i + 1}$ term of $f$. Then, what remains after subtracting the $p_{i + 1}$ term out is a polynomial in $\mathcal P_{i}(\F)$. Because we know that $p_0, \dots, p_i$ is a basis of $\mathcal F_i(\F)$, we know we can produce the rest of the polynomial through a linear combination of $p_0, \dots, p_i$. Thus, $p_0, \dots, p_{i + 1}$ spans $\mathcal P_{i + 1}(\F)$. Because $\mathcal P_{i + 1}(\F)$ has dimension $i + 2$ and we have a list of $i + 2$ vectors that span it, we know that this list is also linearly independent, and thus we know that $p_0, \dots, p_{i + 1}$ is a basis for $\mathcal P_{i + 1}(\F)$. 

                Therefore by induction for any $m \in \mathbb N$, $p_0, \dots, p_m$ is a basis for $\mathcal P_m(\F)$.
            \end{proof}


            \item[10] Suppose $m$ is a positive integer. For $0 \leq k \leq m$, let
            \[
                p_k(x) = x^k(1-x)^{m-k}.
            \]
            Show that $p_0, \ldots, p_m$ is a basis of $\mathcal{P}_m(\mathbb{F})$.


            \begin{proof}
                We can prove this by observing that these polynomials are actually quite similiar to the ones from last exercise. For example, $p_m$ is a polynomial that contains non-zero terms for degree $m$ only. $p_{m-1}$ has non-zero terms for degree terms for degrees $m$ and $m-1$. Thus, we can use the same argument as the previous exercise, using induction, but this time in reverse.
            \end{proof}


            \item[11] Suppose $U$ and $W$ are both four-dimensional subspaces of $\mathbb{C}^6$. Prove that there exist two vectors in $U \cap W$ such that neither of these vectors is a scalar multiple of the other.
            
            \begin{proof}
                The dimension of $U + W$ is at most the dimension of the space itself (6). Therefore, we have that 

                $$6 \geq \dim(U + W) = \dim(U) + \dim(W) - \dim(U \cap W)$$

                $$6 \geq 4 + 4 - \dim(U \cap W)$$

                $$2 \leq \dim(U \cap W)$$

                Now, we can use 2 basis vectors for $\dim(U \cap W)$ as our pair of linearly independent vectors.
            \end{proof}


            \item[15] Suppose $V$ is finite-dimensional and $V_1$, $V_2$, $V_3$ are subspaces of $V$ with $\dim V_1 + \dim V_2 + \dim V_3 > 2 \dim V$. Prove that $V_1 \cap V_2 \cap V_3 \neq \{0\}$.
            
            \begin{proof}
                \begin{align*}
                    \dim V_1 + \dim V_2 + \dim V_3 &> 2 \dim V \\
                    &\geq \dim V + \dim V_1 + \dim V_2 - \dim(V_1 \cap V_2) \\
                    \dim V_3 &> \dim V - \dim(V_1 \cap V_2) \\
                    &\geq \dim(V_1 \cap V_2) + \dim(V_3) - \dim(V_1 \cap V_2 \cap V_3) - \dim(V_1 \cap V_2) \\
                    \dim(V_1 \cap V_2 \cap V_3) &> 0
                \end{align*}
            \end{proof}

            \item[17] Suppose that $V_1, \ldots, V_m$ are finite-dimensional subspaces of $V$. Prove that $V_1 + \cdots + V_m$ is finite-dimensional and 
            \[\dim(V_1 + \cdots + V_m) \leq \dim V_1 + \cdots + \dim V_m.\]
            
            \begin{proof}
                We prove by induction on $m$. For $m = 1$, the statement is trivial. 
                
                Assume the result holds for $m-1$ subspaces. By the inductive hypothesis, $V_1 + \cdots + V_{m-1}$ is finite-dimensional. Then $V_1 + \cdots + V_m = (V_1 + \cdots + V_{m-1}) + V_m$ is a sum of two finite-dimensional subspaces. By the dimension formula for two subspaces,
                \[\dim((V_1 + \cdots + V_{m-1}) + V_m) \leq \dim(V_1 + \cdots + V_{m-1}) + \dim V_m.\]
                
                By the inductive hypothesis, $\dim(V_1 + \cdots + V_{m-1}) \leq \dim V_1 + \cdots + \dim V_{m-1}$. Therefore,
                \[\dim(V_1 + \cdots + V_m) \leq \dim V_1 + \cdots + \dim V_{m-1} + \dim V_m.\]
                
                Thus by induction, the result holds for all positive integers $m$.
            \end{proof}

            \item[19] Explain why you might guess, motivated by analogy with the formula for
the number of elements in the union of three finite sets, that if $V_1$, $V_2$, $V_3$
are subspaces of a finite-dimensional vector space, then
\begin{align*}
\dim(V_1 &+ V_2 + V_3) =  \\
&\dim V_1 + \dim V_2 + \dim V_3 \\
&- \dim(V_1 \cap V_2) - \dim(V_1 \cap V_3) - \dim(V_2 \cap V_3) \\
&+ \dim(V_1 \cap V_2 \cap V_3).
\end{align*}

Then either prove the formula above or give a counterexample.

\begin{proof}
    This statement is FALSE. Consider 3 different lines in $\R^2$. We would have $2 = 1 + 1 + 1 - 0 - 0 - 0 + 0 = 3$!
\end{proof}


\item[20] Prove that if $V_1$, $V_2$, and $V_3$ are subspaces of a finite-dimensional vector space, then
            \begin{align*}
                \dim(V_1 + V_2 + V_3) &= \dim V_1 + \dim V_2 + \dim V_3 \\
                &\quad - \frac{\dim(V_1 \cap V_2) + \dim(V_1 \cap V_3) + \dim(V_2 \cap V_3)}{3} \\
                &\quad + \frac{\dim((V_1 + V_2) \cap V_3) + \dim((V_1 + V_3) \cap V_2) + \dim((V_2 + V_3) \cap V_1)}{3}
            \end{align*}

            \begin{proof}
                This is really 3 formulas squished into one. 

                We have 
                \begin{align*}
                \dim((V_1 + V_2) + V_3) &= \dim(V_1 + V2) + \dim(V_3) - \dim((V_1 + V_2) \cap V_3) \\
                &= \dim(V_1) + \dim(V_2) + \dim(V_3) - \dim(V_1 \cap V_2) - \dim((V_1 + V_2) \cap V_3)
                \end{align*}

                We can apply this again for other permutations to obtain three formulas. Summing them together and dividing by three will give the desired formula. I would write it up but I \emph{really} gotta get to chapter 3 now.
            \end{proof}



\end{enumerate}